% Central resource bundle: edit this file to personalize every report.
% Wszystkie wspólne dane i teksty raportów znajdują się tutaj.

% Dane autora i kursu
\newcommand{\ReportAuthor}{Jakub Jaworski}
\newcommand{\ReportStudentID}{70203}
\newcommand{\ReportGroup}{L2}
\newcommand{\ReportInstitution}{Uniwersytet VIZJA}
\newcommand{\ReportFaculty}{Informatyka dla Inżynierów, Studia Magisterskie}
\newcommand{\ReportCourse}{Zaawansowane zagadnienia algorytmiki i programowania}
\newcommand{\ReportInstructor}{dr hab. inż. Agnieszka Duraj}
\newcommand{\ReportPlace}{Warszawa}
\newcommand{\ReportDate}{\today}

% Stałe elementy strony tytułowej i podpisów
\newcommand{\TextProject}{Projekt 1: Algorytmy grafowe -- analiza algorytmów najkrótszej ścieżki}
\newcommand{\TextAuthor}{Autor}
\newcommand{\TextStudentID}{Nr albumu}
\newcommand{\TextGroup}{Grupa}
\newcommand{\TextInstructor}{Prowadząca}
\newcommand{\TextReport}{Raport}

% Tytuły czterech raportów
\newcommand{\ReportOneTitle}{Podstawy teoretyczne i~plan eksperymentów}
\newcommand{\ReportTwoTitle}{Implementacja podstawowa i~testy poprawności}
\newcommand{\ReportThreeTitle}{Implementacja rozszerzona i~eksperymenty}
% Podział wiersza na okładce zapobiega osamotnionemu słowu „wnioski”.
\newcommand{\ReportFourTitle}{Analiza porównawcza,\\wizualizacja wyników i wnioski}
\newcommand{\ReportOneNumber}{1}
\newcommand{\ReportTwoNumber}{2}
\newcommand{\ReportThreeNumber}{3}
\newcommand{\ReportFourNumber}{4}

% Teksty pomocnicze dla pustych sekcji. Można je usunąć po rozpoczęciu pisania.
\newcommand{\PlaceholderNote}{\textit{Miejsce na treść tej sekcji.}}
\newcommand{\PlaceholderFigure}{\textit{Miejsce na rysunek lub wykres.}}
\newcommand{\PlaceholderTable}{\textit{Miejsce na tabelę wyników.}}
